% Theme: Preprint  --  best for paper summaries / journal-club notes
% Build:  ./build.sh example-preprint.tex   (or compile on Overleaf)
\documentclass[12pt]{beamer}
\usetheme{igpreprint}

\ighandle{@quantum_danesh}
\igseries{Notas de artículos}
\igissue{03}

\begin{document}

\coverblurb{El randomized benchmarking mide la calidad promedio de un
conjunto de puertas. El cycle benchmarking hace una pregunta más precisa:
¿cuál es la fidelidad del ciclo concreto de puertas que de verdad
ejecutas? La responde sin fiarse de la preparación ni de la medición, con
un coste de muestreo que no crece con el número de qubits.}
\coverslide{Artículo comentado}
  {Characterizing large-scale quantum computers via cycle benchmarking}
  {A.~Erhard, J.~J.~Wallman et al.\\ Nature Communications \textbf{10}, 5347 (2019)}

\begin{frame}{La pregunta}{\S1}
El randomized benchmarking (RB) da un único número para \hl{todo un
conjunto de puertas}, promediado sobre secuencias aleatorias de Clifford.

Un algoritmo real, en cambio, ejecuta una y otra vez un \key{ciclo} fijo
$\mathcal{G}$: una capa de puertas en paralelo. Queremos la fidelidad de
proceso de su versión ruidosa
\[
  \tilde{\mathcal{G}} = \mathcal{G}\circ\Lambda ,
\]
medida \hl{sin fiarnos} de la preparación y la medición (SPAM).
\end{frame}

\begin{frame}{Aleatorizar el ruido}{\S2}
Intercalar capas aleatorias de Paulis convierte $\Lambda$ en un canal de
Pauli, que actúa de forma diagonal sobre los operadores de Pauli:
\begin{equation}
  \Lambda(P) = \mhl{\lambda_P}\, P , \qquad P \in \{I,X,Y,Z\}^{\otimes n}.
\end{equation}
La fidelidad de proceso es el promedio de estos autovalores:
\begin{equation}
  F(\tilde{\mathcal{G}},\mathcal{G})
  = \frac{1}{4^n}\sum_{P}\explain{\lambda_P}{cuánto sobrevive $P$}
\end{equation}
Basta con estimar $\lambda_P$ para unos pocos $P$ aleatorios.
\end{frame}

\begin{frame}{El protocolo}{\S3}
\begin{enumerate}
  \item Elige un Pauli aleatorio $P$ y prepara un autoestado suyo de
        autovalor $\pm1$.
  \item Aplica $m$ rondas de \emph{capa aleatoria de Paulis y luego
        $\mathcal{G}$}.
  \item Mide $P$ y anota el valor esperado $f_P(m)$.
  \item Repite con dos profundidades $m_1 < m_2$ y varios $P$ aleatorios.
\end{enumerate}
\begin{notebox}[Esquema]
Para un ciclo de Clifford, $P$ recorre su órbita bajo $\mathcal{G}$, así
que $\lambda_P^{\,m}$ representa un producto de autovalores a lo largo de
esa órbita.
\end{notebox}
\end{frame}

\begin{frame}{Por qué el SPAM desaparece}{\S4}
Cada valor esperado decae geométricamente con la profundidad,
\begin{equation}
  f_P(m) \simeq \mhl<2->{A_P}\,\lambda_P^{\,m},
\end{equation}
\uncover<2->{y todo el error de SPAM queda absorbido en el prefactor $A_P$.}
\uncover<3->{Al dividir dos profundidades, se cancela:
\begin{equation}
  \hat\lambda_P = \left(\frac{f_P(m_2)}{f_P(m_1)}\right)^{\!\mhl*<3->{1/(m_2-m_1)}}.
\end{equation}}
\end{frame}

\begin{frame}{Conclusiones}{\S5}
\begin{itemize}
  \item Mide el \key{ciclo que de verdad ejecutas}, no un promedio sobre
        todo el conjunto de puertas.
  \item Robusto al SPAM: queda en $A_P$ y se cancela en el cociente.
  \item Escalable: los Paulis que hay que muestrear para una precisión
        fija no crecen con $n$.
  \item Supone ruido markoviano que apenas depende de qué capa de Paulis
        se aplicó.
\end{itemize}
\vspace{1mm}
{\footnotesize\color{igMuted}
[1] A.~Erhard et al., Nat.\ Commun.\ \textbf{10}, 5347 (2019).\par
[2] J.~J.~Wallman, J.~Emerson, Phys.\ Rev.\ A \textbf{94}, 052325 (2016).\par}
\end{frame}

\closingslide{Lee la derivación completa}
  {El desarrollo completo, con el álgebra de los autovalores de Pauli,
   está en el blog.}

\end{document}
